\chapter{Linux 本地 CPU 推理引擎项目}

第二册从一个明确的大作业开始：在 Linux 本地实现一个 CPU 版本的量化推理引擎。这个项目不依赖云服务，不调用现成深度学习框架完成核心推理，不把算子实现藏在黑盒库里。它的目标不是一开始就追求工业级完整度，而是让读者亲手打通推理系统的骨架：模型文件、张量、量化权重、算子、线程运行时、推理图、正确性测试和性能报告。

\section{为什么先定义项目}

如果直接从 AI 名词讲起，读者很容易分不清哪些概念服务于训练，哪些概念服务于推理，哪些概念只是框架 API。推理引擎给我们一个清晰边界：输入是一段数据和一个模型文件，输出是模型的预测结果；中间过程包括加载参数、组织张量、执行算子、管理内存、调度线程和记录性能。只要这条路径跑通，后面的矩阵乘、量化、Softmax、Attention 才有落脚点。

本册的教学项目命名为 Linux CPU Quantized Inference，简称 LCQI。它先实现一个量化 MLP 模型，再逐步扩展到更复杂的算子。选择 MLP 作为第一阶段，是因为它包含推理引擎最基本的部件：输入张量、int8 权重、scale、矩阵向量乘、bias、ReLU 和输出分类。模型小，但路径真实。

\begin{keyidea}
贯穿项目的价值不在于模型大，而在于链路完整。一个小模型只要能加载、能推理、能验证、能测量，就已经具备继续扩展到卷积、Attention 和 Transformer 的工程骨架。
\end{keyidea}

\section{项目边界}

LCQI 的第一阶段只支持本地 Linux 机器上的 CPU 执行，C++20 编写。这里的“本地”不是云端服务，也不是把推理委托给深度学习框架；这里的“Linux”是实验和性能分析的标准平台。模型格式使用教学文本文件，避免读者一开始被复杂二进制格式拖住。权重使用 int8 存储，每一层有一个 scale，把 int8 权重反量化到 float 参与计算。输入和激活先用 float，后续章节再讨论激活量化和更高效的 int8 累加。

第一阶段模型结构固定为两层 MLP：

\begin{itemize}
\item 输入向量长度为 4。
\item 第一层输出隐藏向量，执行量化权重矩阵乘、bias 和 ReLU。
\item 第二层输出分类 logits。
\item 推理结果取最大 logit 的下标。
\end{itemize}

这个结构足够小，可以手工检查结果；也足够真实，因为大模型中的 MLP、投影层和分类头都可以看作更大规模的矩阵运算。

\section{工程模块}

LCQI 至少包含七个模块。第一是模型加载器，把本地模型文件解析成内存中的权重和 bias。第二是张量结构，表达 shape、元素数量和连续数据。第三是量化权重结构，保存 int8 数据、scale、输入维度和输出维度。第四是算子内核，先实现量化线性层和 ReLU。第五是推理执行器，把多层算子按顺序执行。第六是测试，用固定模型和输入验证输出。第七是 benchmark，记录延迟和吞吐。

这些模块会随着章节推进不断变强。讲缓存时，我们检查张量连续性；讲线程时，我们给线性层增加并行输出切分；讲 SIMD 时，我们分析内层点积；讲 AI 基础时，我们把线性层、激活和 logits 放进模型语义；讲量化时，我们改进权重格式和误差检查。

\begin{lstlisting}[language=C++,caption={推理引擎第一阶段的核心对象关系}]
struct QuantizedLinearLayer {
    std::int32_t input_size;
    std::int32_t output_size;
    float scale;
    std::vector<std::int8_t> weights;
    std::vector<float> bias;
};

struct TinyMlpModel {
    QuantizedLinearLayer hidden;
    QuantizedLinearLayer output;
};
\end{lstlisting}

这段结构不是最终设计，但已经表达了推理引擎的基本事实：模型不是魔法对象，而是一组带形状、类型和量化元数据的数组。算子开发的第一步，就是让这些数组的含义清楚。

\section{Linux 环境要求}

本册实验以 Linux 为目标环境。最低要求是能使用 CMake、支持 C++20 的编译器和基本命令行工具。推荐工具包括 \cmd{cmake}、\cmd{ctest}、\cmd{perf}、\cmd{taskset}、\cmd{lscpu}、\cmd{numactl}。手机 Termux 环境可以完成大部分 C++ 构建和小模型推理，但 perf、NUMA 和硬件计数器可能受权限或内核限制。

实验报告必须记录系统信息。至少包括 CPU 型号、核心数、线程数、内存大小、编译器版本、编译选项、操作系统版本和是否运行在容器或手机环境中。性能工程的第一条纪律是：没有环境上下文的数字不能比较。

\begin{labbox}
本章实验：构建 LCQI 项目，运行内置测试，确认模型文件能被加载，推理输出能通过断言。随后记录一次本地环境报告，包括 \cmd{uname -a}、\cmd{lscpu} 和编译器版本。如果某些命令在 Termux 中不可用，要在报告中写明。
\end{labbox}

\section{从本章开始积累证据}

后续每章都要回到这个项目。现代 CPU 结构解释为什么内存布局重要；并发章节解释为什么线程池不能随便写；AI 基础章节解释模型文件中的权重和算子语义；算子章节解释矩阵乘和归一化如何实现；量化章节解释 scale、误差和低精度累加；最后的 benchmark suite 把所有证据汇总成工程报告。

这条路线会慢一点，但它避免了一个常见问题：懂一点 AI 名词，懂一点 SIMD 指令，却不知道怎样把它们接成能运行的系统。本册的目标是把这条链接起来。
